We now describe a general methodology for training a model that can in-context learn a function class $\mathcal{F}$ with respect to a distribution $D_{\mathcal{F}}$ over functions, and $D_\mathcal{X}$ over inputs. To do so, we start by constructing random training prompts as follows.
For each prompt, we first sample a random function $f$ from the class according to $D_{\mathcal{F}}$, then create a set of random inputs $x_1, \ldots, x_{k+1}$ drawn independently from  $D_{\mathcal{X}}$, and finally evaluate $f$ on these inputs to produce the prompt
$P = (x_1, f(x_1), \ldots, x_{k+1}, f(x_{k+1}))$.
For example, in the case of linear functions, inputs could be drawn from the isotropic Gaussian distribution $N(0, I_d)$, and a random function chosen by sampling weight vector $w$ from $N(0, I_d)$ and setting $f(x) = w^\top x$.

Now, given such prompts, we train a model to predict $f(x_i)$ for a given $x_i$ based on a set of preceding in-context examples.
Concretely, let $P^{i}$ denote the prompt prefix containing $i$ in-context examples (the first $i$ input-output pairs) and the $(i+1)^{\text{th}}$ input: $P^i = (x_1, f(x_1), x_2, f(x_2), \ldots, x_{i}, f(x_{i}), x_{i+1})$. 
Then, we train a model $M_\theta$ parameterized by $\theta$ aiming to minimize the expected loss over all the prompt prefixes:
\begin{equation}
\min_\theta \ 
\mathds{E}_{P}\left[\frac{1}{k+1}
\sum_{i=0}^k \ell\left(M_\theta\left(P^i\right), f\left(x_{i+1}\right)\right) 
\right],
\label{eq:training_objective}
\end{equation}
where $\ell(\cdot, \cdot)$ is an appropriately chosen loss function.
Below, we describe how this general methodology can be implemented for a concrete model family (see Appendix~\ref{app:setup} for additional details).

\paragraph{Model structure.}
We use a decoder-only Transformer architecture~\citep{vaswani2017attention} from the GPT-2 family~\citep{radford2019language}. 
Our model consists of 12 layers, 8 attention heads, and a 256-dimensional embedding space (9.5M parameters).
This architecture takes as input a sequence of vectors in its embedding space and predicts the next vector in the sequence within the same space (in language modeling, these vectors correspond to input tokens).
We apply this architecture to our prompt format of $(x_1,f(x_1),\ldots,x_{k+1},f(x_{k+1}))$ as follows.
We map each prompt output $f(x_i)$ to the same dimension as prompt inputs $x_i$ by appending zeros, and map the prompt inputs and outputs into the latent embedding space of the Transformer through a (learnable) linear transformation.
We then use another (learnable) linear transformation to map the vector predicted by the model to a scalar.
Note that the Transformer architecture allows us to compute the prediction ($M_\theta(P^i)$) for all prompt prefixes in a single forward pass.

\paragraph{Training.}
We train the model according to the training objective in~\eqref{eq:training_objective} using squared error as the loss function. 
We do so by sampling a batch of random prompts at each training step and then updating the model through a gradient update (we use a batch size of 64 and train for 500k total steps). 
This training is done from scratch, that is, we do \emph{not} fine-tune a pre-trained language model, nor do we train on actual text.

\paragraph{Curriculum learning.}
Many natural function classes contain functions of varying complexity. We exploit this by training our model using a curriculum \citep{bengio2009curriculum, elman1993learning, sanger1994neural, wu2020curricula}, where we train on a simpler distribution of functions in the beginning (e.g., linear functions with weight vectors restricted to a low-dimensional subspace) and gradually increase the function complexity.
This speeds up training drastically, 
often allowing us to train models that would be significantly more expensive to train without a curriculum (see Section \ref{sec:factors} for details).